\scanpage{4}
% archive foliation: 3
\diagram{\begin{tikzcd}[column sep=small]
0\arrow[r]&\Omega\arrow[r]&E\arrow[r,"p"]&\mathcal{O}_S\arrow[r]&0
\end{tikzcd}}

$P=\mathbb{P}(E)$, $\sigma$ section définie par $p$, $\widehat P$ complété formel de $P$ le long de $\sigma$
\[
\widehat S=\mathcal{O}_{\widehat P}
\]
$\Omega$ opère \inserted{à g.} sur $\widehat S$ : car il opère sur l'extension (à g.) donc sur $(P,\sigma)$ et $\widehat P$ (à droite)

$L=\mathcal{O}_P(1)$, $\widehat L$ module inv. sur $\widehat S$

$\Omega$ opère à g. aussi sur $\widehat L$, de façon compatible avec opérations de $\Omega$ sur l'Anneau de scalaires $\widehat S$

Cas $E$ \inserted{extension} splittée (i.e. on a donné splitting $s:\mathcal{O}_S\to E$)
\[
E^0\simeq\Omega\times\mathcal{O}_S
\]
Alors \deleted{\unsure{$\widehat P(E^0)$} \unsure{par l'anneau comme} \illegible{2}},

\deleted{\illegible{1}} $\widehat P(S')\simeq$ Quotients inversibles \inserted{$L'$} de $E$, qui sur $S$ réd. coïncident avec le quotient donné $\mathcal{O}_S$ --- donc l'image dans $L'$ de la section $e$ de $\mathcal{O}_S\subset E$ est une base de $L'$, permettant d'identifier $L'$ à $\mathcal{O}_S$, de sorte que le quotient $\unsure{E}\to\deleted{\unsure{L'}}\inserted{\mathcal{O}_S}$ est défini par $(x,\lambda)\longmapsto\varphi(x)+\lambda$, où $\varphi:\Omega\to\mathcal{O}_S$ hom bien défini, de $\Omega$ dans nilradical \inserted{$N_S$}.

Donc
\[
\begin{aligned}
\widehat P^0&=\mathrm{V}(\Omega)^\wedge\quad\text{(complété formel le long de la section nulle)}\\
&=\operatorname{Spf}\underbrace{\widehat{\operatorname{Sym}}(\Omega)}_{\widehat S^0}
\qquad \widehat L^0\simeq\widehat S^0
\end{aligned}
\]
Il faut expliciter l'op. de $\Omega$ sur $\widehat S$. Soit $a\in\Omega(S)$.

Donc $a$ opère sur l'ext. $E$ par
\[
u_a(x,\lambda)=(x+\lambda a,\lambda)
\]
et si \deleted{\unsure{Donnons}} $\varphi:\Omega\to N_S$, définissant le \inserted{quotient} $\widetilde\varphi$ de $E$, \inserted{$\widetilde\varphi:E\xrightarrow{\text{épi}}\mathcal{O}_S$}, par $\widetilde\varphi(x,\lambda)=\varphi(x)+\lambda$, on aura le nouveau quotient par
\[
\begin{aligned}
\widetilde\varphi\circ u_a:\ (x,\lambda)&\longrightarrow\varphi(x+\lambda a)+\lambda\\
&=\varphi(x)+(1+\varphi(a))\lambda\\
&=(1+\varphi(a))\left(\frac{\varphi\deleted{\unsure{(x)}}}{1+\varphi(a)}(x)+\lambda\right),
\end{aligned}
\]
c'est aussi le quotient défini par \inserted{$(x,\lambda)\mapsto$} $\varphi^a(x)+\lambda$, donc
